\section{Finite distribution geometry}
\label{app:finite-distribution-geometry}

This appendix records finite-space consequences of the distribution-relative receiver Bayes risk that are useful for structural selection analysis but not needed for the main theorem spine.

Assume throughout that $X$, $Y$, and all receiver images are finite and specialize to
\[
\mathsf A=Y,
\qquad
\ell(y,\hat y)=\mathbf 1[y\neq\hat y].
\]

\subsection{Threshold stability away from a boundary}

For a tolerance $\varepsilon\in[0,1]$, define
\[
\gamma_{\mu,\varepsilon}
:=
\min_{K\subseteq J}
|\delta_\mu(K)-\varepsilon|.
\]

\begin{corollary}[Localization-family stability radius]
\label{cor:app-localization-family-stability}
If $\gamma_{\mu,\varepsilon}>0$ and
\[
d_{\TV}(\mu,\nu)<\gamma_{\mu,\varepsilon},
\]
then
\[
\Rcal_{\mu,\varepsilon}
=
\Rcal_{\nu,\varepsilon}.
\]
Hence their inclusion-minimal localization antichains are also equal.
\end{corollary}

\begin{proof}
The main-text TV-Lipschitz theorem implies
\[
|\delta_\mu(K)-\delta_\nu(K)|
<\gamma_{\mu,\varepsilon}
\]
for every $K$. Every defect value therefore remains on the same side of the threshold $\varepsilon$.
\end{proof}

\subsection{Finite receiver risk is piecewise linear}

For fixed $K$, write
\[
p_{z,y}^\mu
:=
\mu\{x:Q_K(x)=z,\ \Phi(x)=y\}.
\]

\begin{proposition}[Finite receiver risk is ordinary Bayes classification error]
\label{thm:app-bayes-defect-formula}
For every receiver family $K$,
\[
\delta_\mu(K)
=
1-\sum_{z\in Z_K}\max_{y\in Y}p_{z,y}^\mu.
\]
In particular, $\mu\mapsto\delta_\mu(K)$ is concave and piecewise linear on the probability simplex $\Delta(X)$.
\end{proposition}

\begin{proof}
For a rule $h:Z_K\to Y$,
\[
R_\mu(K,h)
=
1-\sum_{z\in Z_K}p_{z,h(z)}^\mu.
\]
Each receiver value can be optimized independently, so the largest correct mass at $z$ is $\max_y p_{z,y}^\mu$.
\end{proof}

\subsection{Polyhedral selection geometry}

\begin{proposition}[Finite localization cells]
\label{thm:app-polyhedral-localization-cells}
Fix a finite receiver family $J$ and tolerance $\varepsilon\in[0,1]$. There is a finite polyhedral subdivision of $\Delta(X)$ such that on the relative interior of every cell:
\begin{enumerate}[label=(\roman*)]
\item the optimizer set $\Opt_\mu(K)$ is constant for every $K\subseteq J$;
\item the thresholded localization family $\Rcal_{\mu,\varepsilon}$ is constant;
\item the inclusion-minimal localization antichain
\[
\Mcal_{\mu,\varepsilon}:=\Min_{\subseteq}\Rcal_{\mu,\varepsilon}
\]
is constant.
\end{enumerate}
\end{proposition}

\begin{proof}
For fixed $K$ and $h$, the map
\[
\mu\longmapsto R_\mu(K,h)
\]
is affine. Partition the finite decoder set into classes whose members induce identical affine risk functionals. Distinct classes can exchange optimizer status only on nontrivial tie hyperplanes
\[
R_\mu(K,h)=R_\mu(K,h'),
\]
and threshold membership can change only on nontrivial hyperplanes
\[
R_\mu(K,h)=\varepsilon.
\]
The finite arrangement of these hyperplanes together with the faces of $\Delta(X)$ yields a finite polyhedral subdivision. On each relative interior the ordering and equality pattern of all distinct risk functionals, and their signs relative to $\varepsilon$, are fixed. The optimizing equivalence classes are therefore fixed; every decoder inside one such class enters or leaves the optimizer set together, so the full set $\Opt_\mu(K)$ is fixed as well. The thresholded family and its inclusion-minimal members are then fixed on the same cell.
\end{proof}

We use \emph{localization cell} for a cell of such a subdivision and reserve \emph{input regime} or \emph{context} for the scientific population being analyzed. The distinction avoids overloading the word ``regime.''

The main text uses only the consequence needed for interpretation: quantitative receiver risks vary continuously under TV shift, while optimizer identities and thresholded selections can change discretely at decision boundaries.
